% =====================================================================
%  UniMe Beamer template (16:9)
%  University of Messina, Department of Engineering.
%
%  Files:   template.tex          this file: copy and rename it
%           beamerthemeunime.sty  the theme (colours, title page, footline,
%                                 section dividers, closing slide)
%           assets/               UniMe seal, original colours and white
%           figures/              put your images here
%
%  Build:   latexmk -pdf template.tex      (or pdflatex twice)
%  Overleaf: zip the folder and use "New Project > Upload Project".
%  Docs:    README.md
% =====================================================================
\documentclass[aspectratio=169,10pt]{beamer}
\usetheme{unime}

\usepackage{amsmath,amssymb}
\usepackage{booktabs,tabularx,multirow,array}
\usepackage{pifont}      % \ding{51} check, \ding{55} cross
\usepackage{listings}    % code listings
\usepackage{qrcode}      % QR code on the closing slide
\usefonttheme[onlymath]{serif}
\graphicspath{{figures/}}

% handy marks and column types
\newcommand{\yes}{\textcolor{unimegreen!70!black}{\ding{51}}}
\newcommand{\no}{\textcolor{unimered}{\ding{55}}}
\newcolumntype{L}[1]{>{\raggedright\arraybackslash}p{#1}}
\newcolumntype{Y}{>{\raggedright\arraybackslash}X}
\lstset{basicstyle=\ttfamily\scriptsize,columns=fullflexible,keepspaces=true,
  backgroundcolor=\color{unimelight},frame=single,rulecolor=\color{unimeblue!40},showstringspaces=false}

% ---------------------------------------------------------------------
% Title page data. The optional [short] arguments fill the footline.
% ---------------------------------------------------------------------
\title[Short title]{Title of the Presentation}
\subtitle{Subtitle or one-sentence message}
\author[N.~Surname et al.]{Name Surname, Second Author, Third Author}
\institute{University of Messina \\ Department of Engineering}
\date{}  % leave empty to hide the date

% Theme colours available in the document:
%   unimeblue (main), unimedark, unimegrey, unimelight (block backgrounds),
%   unimeink (text), unimeorange (alerts), unimegreen, unimered

\begin{document}

\begin{frame}[plain,noframenumbering]
  \titlepage
\end{frame}

% Every \section opens with a blue outline slide that highlights it.
\section{Introduction}

\begin{frame}{A simple slide}
\framesubtitle{Frame subtitle, shown in grey under the title}
\begin{itemize}
  \item First point, with an \alert{alerted} word
  \item Second point
  \begin{itemize}
    \item A sub-point
  \end{itemize}
  \item Third point
\end{itemize}
\end{frame}

\begin{frame}{Two columns: text and image}
\begin{columns}[c]
\column{0.5\textwidth}
\begin{itemize}
  \item Text on the left
  \item An image on the right, from \texttt{figures/}
  \item Replace \texttt{example-image-a} with your file
\end{itemize}
\column{0.46\textwidth}
\includegraphics[width=\linewidth]{example-image-a}
\end{columns}
\end{frame}

\begin{frame}{Blocks}
\begin{columns}[t]
\column{0.32\textwidth}
\begin{block}{Block}
Standard block, UniMe blue.
\end{block}
\column{0.32\textwidth}
\begin{exampleblock}{Example block}
Green block, for results or definitions.
\end{exampleblock}
\column{0.32\textwidth}
\begin{alertblock}{Alert block}
Orange block, for warnings or limits.
\end{alertblock}
\end{columns}
\vspace{0.4cm}
\begin{center}
\begin{tikzpicture}
  \node[draw=unimeblue,line width=1.2pt,fill=unimelight,rounded corners=6pt,text width=0.75\textwidth,align=center,inner sep=10pt,font=\large]
  {A key message in a highlighted box.};
\end{tikzpicture}
\end{center}
\end{frame}

\section{Method}

\begin{frame}{Equations}
\begin{block}{A displayed equation}
\[ \Pr[\text{event}] \;=\; 1-e^{-\lambda\Delta} \;\approx\; \lambda\Delta \quad \text{for small } \lambda\Delta \]
\end{block}
Inline math such as ${a(t)\le d}$ can be kept on one line with braces.
\end{frame}

\begin{frame}{A table}
\begin{center}
\footnotesize
\renewcommand{\arraystretch}{1.2}
\begin{tabularx}{0.85\textwidth}{@{}lYcc@{}}
\toprule
\textbf{Method} & \textbf{Description} & \textbf{Property A} & \textbf{Property B} \\
\midrule
Baseline & a short description & \no & \no \\
Variant & another description & \yes & \no \\
Proposed & our approach & \yes & \yes \\
\bottomrule
\end{tabularx}
\end{center}
\end{frame}

\begin{frame}[fragile]{Code}
\begin{lstlisting}
{ "id": "example",
  "payload": { "action_type": "RESUME", "target_entity": "line_1" } }
\end{lstlisting}
\end{frame}

\section{Results}

\begin{frame}{A full-width figure}
\begin{center}
\includegraphics[width=0.7\linewidth,height=0.6\textheight,keepaspectratio]{example-image-b}
\end{center}
\begin{itemize}
  \item One sentence that states what the figure shows
\end{itemize}
\end{frame}

\section{Conclusions}

\begin{frame}{Conclusions}
\begin{enumerate}
  \item First take-away
  \item Second take-away
  \item Third take-away
\end{enumerate}
\vspace{0.2cm}
\textbf{Next}: future work in one line.
\end{frame}

% Closing slide: blue, with the title and "Thank you"; the argument is
% printed below (contacts, link, QR code). Leave it empty with \closingframe{}.
\closingframe{%
  \begin{tabular}{@{}c@{\hspace{0.8cm}}l@{}}
  \raisebox{-0.5\height}{\colorbox{white}{\color{black}\qrcode[height=1.5cm]{https://www.unime.it}}} &
  \begin{tabular}{@{}l@{}}Name Surname --- \texttt{name.surname@unime.it}\\ \texttt{https://www.unime.it}\end{tabular}
  \end{tabular}}

\end{document}
